\documentclass[12pt]{article}

% Idioma y codificación
\usepackage[spanish]{babel}
\usepackage[utf8]{inputenc}

% Márgenes
\usepackage{geometry}
\geometry{margin=2.5cm}

% Paquetes útiles
\usepackage{graphicx} % Imágenes
\usepackage{listings} % Código
\usepackage{xcolor}

% Configuración para código en C
\lstset{
    language=C,
    basicstyle=\ttfamily\small,
    numbers=left,
    numberstyle=\tiny,
    frame=single,
    breaklines=true
}

\begin{document}

% ---------- PORTADA ----------
\begin{titlepage}
    \centering
    \vspace*{4cm}
    
    {\Large \textbf{Programación Estructurada}}\\[1cm]
    {\large Práctica No. 4}\\[2cm]
    
    {\large Alexis Roberto Constantino Núñez}\\

    {\large Matricula: 2201619}\\
    {\large Grupo 531}\\[1cm]
    
    {\large \today}
\end{titlepage}

% ---------- PROBLEMA 1 ----------
\section{Problema 1}

\subsection{Enunciado del problema}
%Aquí va una breve explicación del programa realizado en lenguaje C.  
Dise\~{n}e e implemente un programa basado en p\'{a}rametros por referencia que convierta una cantidad en pesos a yenes, d\'{o}lares, libras esterlinas y euros.
\subsection{Especificación de las entradas y salidas}
Entradas: Cantidad de pesos a convertir.\\
Salidas: Cantidad equivalente a la conversi\'{o}n.
\subsection{Algoritmo del programa}
\begin{verbatim}
PREGUNTAR a que moneda quiere convertir
PREGUNTAR la cantidad de pesos a convertir
CALCULAR la conversion deacuerdo a la opcion y la cantidad de pesos
IMPRIMIR el resultado de la conversion
\end{verbatim}
\subsection{Pruebas de escritorio }
\begin{verbatim}
[1] Conversor de moneda
[2] Encuesta de vacunacion
[3] Calculo de finiquito
[4] Salir
1
Conversion de:
[1] Pesos a Yenes
[2] Pesos a Dolares
[3] Pesos a Libras Esterlinas
[4] Pesos a Euros
1
Ingrese pesos
17.5
Res: 159.0909
Conversion de:
[1] Pesos a Yenes
[2] Pesos a Dolares
[3] Pesos a Libras Esterlinas
[4] Pesos a Euros
2
Ingrese pesos
17.9
Res: 1.0176
Conversion de:
[1] Pesos a Yenes
[2] Pesos a Dolares
[3] Pesos a Libras Esterlinas
[4] Pesos a Euros
3
Ingrese pesos
19
Res: 0.8075
Conversion de:
[1] Pesos a Yenes
[2] Pesos a Dolares
[3] Pesos a Libras Esterlinas
[4] Pesos a Euros
4
Ingrese pesos
43.2
Res: 2.1094
[1] Conversor de moneda
[2] Encuesta de vacunacion
[3] Calculo de finiquito
[4] Salir
4
Saliendo...
\end{verbatim}
% ---------- PROBLEMA 2 ----------
\section{Problema 2}

\subsection{Enunciado del problema}
%Aquí va una breve explicación del programa realizado en lenguaje C.  
En el cruce peatonal de una frontera se realiz\'{o} una encuesta sobre vacunaci\'{o}n a 500 personas y se les pregunt\'{o} sobre vacunaci\'{o}n. Dise\~{n}e una funci\'{o}n mediante un ciclo que lea las respuestas de las encuestas y regrese mediante par\'{a}metros por referencia la cantidad de vacunas aplicadas de cada tipo 1. J&J, 2. Cansino, #. AstraZeneca, 4.Pfizer 5.No se vacun\'{o}.
\subsection{Especificación de las entradas y salidas}
Entradas: Ninguna, la lista frue creada artificialmente.\\
Salidas: El conteo correspondiente de la lista.
\subsection{Algoritmo del programa}
\begin{verbatim}
CREAR lista con N cantidad de valores
IMPRIMIR el conteo de la lista creada
\end{verbatim}
\subsection{Pruebas de escritorio }
\begin{verbatim}
Encuesta de Vacunacion a N personas

Datos   Cantidad
 --------------
   1    1
   2    1
   3    2
   4    2
   5    3
 
\end{verbatim}
%-----------------Problema 3---------------------
\section{Problema 3}

\subsection{Enunciado del problema}
El finiquito corresponde al pago de las prestaciones que se establecen en la ley Federal de Trabajo, se otorga al presentarse una renuncia por parte del trabajador. Las prestaciones son: 

A. Aguinaldo: Corresponden 15 d\'{i}as por a\~{n}o trabajado, para calcular la parte proporcional por d\'{i}a se divide 15/365 y eso se multiplica por los d\'{i}as trabajados y el sueldo diario.

B. Vacaciones: Corresponden 6 d\'{i}as por a\~{n}o trabajado, para calcular la parte proporcional por d\'{i}a se divide 6/365 y eso se multiplica por los d\'{i}as trabajados, esto representa los d\'{i}as que le corresponden de vacaciones. La prima vacacional se obtiene multiplicando los d\'{i}as de vacaciones por el sueldo diario y calculando el 25 porciento del resultado anterior.
\subsection{Especificacion de las entradas y salidas}
Entradas: Salario diario y Dias trabajados\\ 
Salidas: Dias de aguinaldo y de vacaciones, Cantidad moneraria del aguinaldo y vacaciones, Finiquito.
\subsection{Algoritmo del programa}
\begin{verbatim}

\end{verbatim}
\subsection{Pruebas de escritorio}
\begin{verbatim}
PREGUNTAR sueldo diario
PREGUNTAR dias trabajados
CALCULAR dias de aguinaldo: dias trabajados * 15/365
CALCULAR aguinaldo: dias de aguinaldo * sueldo diario
CALCULAR dias de vacaciones: dias trabajados * 6/365
CALCULAR prima vacacional: dias de vacaciones * sueldo * 0.25
CALCULAR finiquito: aguinaldo + primavacacional
IMPRIMIR dias aguinaldo y aguinaldo
IMPRIMIR dias de vacaciones y primavacacional
IMPRIMIR finiquito
\end{verbatim}
%-----------------Implementacion en C---------------------
\subsection{Implementacion del algoritmo en C}
\begin{lstlisting}
#include <stdio.h>
//prototipos
/*Problema 1*/
void conversion (void);
void conversor (int opc, float *input);
/*Problema 2*/
void encuesta (void);
void conteo (int datos[], int opciones[]);
/*Problema 3*/
void pago (void);
void calAguinaldo (float dias, float sueldo, float *diasagui, float *aguinaldo);
void calPrimaVac (float dias, float sueldo, float *diasvac, float *primavac);
//main
int main(void)
{
        int opc=0;
        do{
                printf("[1] Conversor de moneda\n[2] Encuesta de vacunacion\n[3] Calculo de finiquito\n[4] Salir\n");//menu
                scanf(" %d", &opc);
                switch(opc)
                {
                        case 1:
                                printf("Conversion de:\n");
                                conversion();
                                break;
                        case 2:
                                printf("Encuesta de Vacunacion a N personas\n");
                                encuesta();
                                break;
                        case 3:
                                printf("Calculo de finiquito\n");
                                pago();
                                break;
                        case 4:
                                printf("Saliendo...\n");
                                break;
                        default: printf("Invalid section\n");
                }
        }while(opc != 4);
        return 0;
}
//Problema 1: Conversion de monedas
void conversion (void)
{
        int opc;
        float input=0;

        printf("[1] Pesos a Yenes\n[2] Pesos a Dolares\n[3] Pesos a Libras Esterlinas\n[4] Pesos a Euros\n");
        scanf("%d", &opc);//opciones
        printf("Ingrese pesos\n"); scanf(" %f", &input);
        conversor(opc, &input);
        printf("Res: %.4f\n", input);
}
void conversor (int opc, float *input)
{
        const float YEN=0.11;
        const float DLL=17.59;
        const float LBE=23.53;
        const float EUR=20.48;
        switch(opc)
        {
                        case 1:
                                *input = *input / YEN;
                                break;
                        case 2:
                                *input = *input / DLL;
                                break;
                        case 3:
                                *input = *input / LBE;
                                break;
                        case 4:
                                *input = *input / EUR;
                                break;
        }
}
//Problema 2: Encuesta de vacunacion
#define PERSONAS 10
#define OPC 5
void encuesta (void)
{
        int datos[PERSONAS] = {5, 3 ,4, 5, 5, 1, 4, 3, 2}; // posibles valores: {1, 2, 3, 4, 5}
        int opciones[OPC]={0};

        conteo (datos, opciones);

        printf("Datos\tCantidad\n --------------\n");
        for(int i=0; i<OPC; i++)
        {
                printf("   %d\t%d\n", i+1, opciones[i]);
        }
}
void conteo (int datos[], int opciones[])
{
        for(int i=0; i<PERSONAS; i++)
        {
                switch (datos[i])
                {
                        case 1: opciones[0] ++; break;
                        case 2: opciones[1] ++; break;
                        case 3: opciones[2] ++; break;
                        case 4: opciones[3] ++; break;
                        case 5: opciones[4] ++; break;
                        default: printf("Opcion Invalida.\n");
                }
        }

}
//Problema 3: Calculo de finiquito
void pago (void)
{
        float dias;
        float sueldo, diasvac, diasagui;
        float aguinaldo, primavac, finiquito;
        aguinaldo=primavac=finiquito=diasvac=0;
        printf("Ingrese sueldo diario:\n"); scanf(" %f", &sueldo);//sueldo
        printf("Ingrese dias trabajados\n"); scanf(" %f", &dias);//dias
        calAguinaldo(dias, sueldo, &diasagui, &aguinaldo);
        calPrimaVac (dias, sueldo, &diasvac, &primavac);
        printf("Aguinaldo:\t %.4f dias.\t=\t%.4f pesos.\n", diasagui, aguinaldo);
        printf("Vacaciones:\t %.4f dias.\t=\t%.4f pesos.\n", diasvac, primavac);
        finiquito = aguinaldo + primavac;
        printf("Finiquito: %.4f\n", finiquito);
}
void calAguinaldo (float dias, float sueldo, float *diasagui, float *aguinaldo)
{
        *diasagui = dias * 15/365;
        *aguinaldo = *diasagui * sueldo;
}
void calPrimaVac (float dias, float sueldo, float *diasvac, float *primavac)
{
        *diasvac = dias * 6/365;
        *primavac = *diasvac * sueldo;//(*diasvac * sueldo) * 0.25;
}
\end{lstlisting}



\end{document}

